\documentclass{article}

\usepackage{amsmath,amssymb}
\usepackage{bbm}

\begin{document}

\section{Piecewise constant distance-based}

Consider smoothing kernels of the form
\begin{align*}
K^\theta(x, x') &{}= k^\theta(d(x, x')) \\
k^\theta(z) &{}= \sum_{j=1}^m \theta_i \mathbbm{1}\{z \le thresh_i\}
\end{align*}
first for a pre-specified $m$ and thresh settings.
......


\begin{align*}
  & \text{arg\,max}_\theta \sum_{i=1}^n \log\left( \frac1{n-1} \sum_{i'=1,i'\neq i}^n K(x_i, x_{i'}) \right)
  \\ & \text{s.t.} \sum_{i=1}^m \theta_m thresh_i = 1
\end{align*}
still have abuse of notation w/ $x$ here vs. above...

plus latex cleanup

\begin{align*}
  & \text{arg\,max}_\theta \sum_{i=1}^n \log\left( \frac1{n-1} \sum_{i'=1,i'\neq i}^n \sum_{j=1}^m \theta_i \mathbbm{1}\{d(x_i, x_{i'}) \le thresh_i\} \right)
  \\ & \text{s.t.} \sum_{i=1}^m \theta_m thresh_i = 1
\end{align*}



see orgzly notes


\section{from dropbox}


\section{Kernel ensemble/shaping}

(Shaping here refers to choosing a kernel function shape, not kernel bandwidth matrix.)

\subsection{Kernel density estimation}

\begin{itemize}
\item 
Potential 1-D bandwidth selection goal:
\[\max_h \sum_{i\in I} \log \left[ \frac{1}{\left|L_i\right| h}
    \sum_{l\in L_i} K\left(\frac{x_i-x_l}{h}\right) \right] \]
Common approach: grid search.
\item
Corresponding 1-D kernel ensemble/shaping goal:
\[\max_{\gamma\in\Delta^{\left|M\right|}} \sum_{i\in I} \log \sum_{m\in M} \gamma_m \frac{1}{\left|L_i\right| h_m}
    \sum_{l\in L_i} K\left(\frac{x_i-x_l}{h_m}\right) \]
\end{itemize}
This is a concave maximization problem.
%
\begin{itemize}
\item

  If $K$ is decreasing with increasing magnitude of its argument, then so does
  the ensemble.

\item

  If $K$ is the boxcar kernel, this is very directly choosing the shape of the
  kernel / values of $K$.
  %
  We could imagine optimizing over all bandwidths instead (optimizing
  $\gamma(h)$), and this may simplify to optimizing over a $\gamma$ vector with
  $H=\{x_i-x_l:i\in I, l\in L_i\}$.
  %
  In fact, this might even correspond to optimizing over all $K$ that are
  decreasing with increasing magnitude arguments.
  %
  However, $\left|I\right|^2$ bandwidths is too many --- we'll need to cut some off.
  %
  Need thought about interpolation and extrapolation past last bandwidth.

\item

  Kernel shaping +
  \href{http://www.stat.columbia.edu/~bodhi/Talks/NPCvxReg.pdf}{convex
    regression}? Can we optimize over all log-concave kernels by specifying
  tangent lines/planes to each encountered $x_i-x_l$ and constraining them?
\item

  Placing $\left|M\right|$ equally-weighted bandwidths? But 1-bandwidth
  selection is an instance of this and is(?) hard\ldots What about for specific
  kernel?

\item

  Something more rigid: kernel bandwidth distribution from a parametric family

  \[\max_\theta \sum_{i\in I} \log \sum_{m\in M} \int_h f(h;\theta) \frac{1}{\left|L_i\right| h}
    \sum_{l\in L_i} K\left(\frac{x_i-x_l}{h} \right) \,\mathrm{d}h \]

  Here, $h$ is an integration variable, not an optimization variable; $\theta$
  is the only optimization variable.
  %
  In fact, we could perhaps replace $h$ with $g(r)$ with $r$ the integration
  variable + potentially adjust according to substitution rule (should just be
  based on $(g^{-1})'(r)$, constant wrt optimization vars --- or may not need at
  all if integration is over $r$).
  %
  If $f$ is concave (log-concave by itself is not sufficient), then this is
  concave maximization. TODO inspect --- may want transformation if
  possible so can cover an unbounded domain).
  %
  If $f$ is convex (log-convex would imply but isn't necessary), then maybe a
  variational-ish alternative could be used.

  Candidate for variationalish approach: exponential distribution, natural
  without need for transformation.  No, not convex in $\lambda$ parm. 

  Candidate for variationalish approach?: uniform over $(0,r)$ and use
  $h=\frac1{1-r}-1$?  There is probably a problem with the discontinuity.

  Candidate for variationalish approach: exponential distribution, using
  $\beta$ parm --- no, still just log-concave, not convex/concave.

  Candidate for nonvariationalish approach: exponential distribution, using
  $\tilde\lambdas$ or $\tilde\beta$ parm? clearly log-concave, but need
  concave\ldots and it's not.

  TODO can marginalization / partial integration (?) log-concavity preservation
  be useful to develop possibilities?

  TODO for variationalish approach?: gamma/beta distributions in non-log-concave
  ($\leadsto$ log-convex???) portions + transform... no, lack of log-concavity
  (at least for beta) is just for the var, not the parms\ldots but can conjugate
  distributions be used for anything here?

  TODO reparameterization of beta distribution based on log moments involving
  digamma function and old parameters?

  for concave: piecewise linear (fixeda,0) -- (varymid,fixed?val) -- (fixedb,0)
  parameterized by slope of one segment (or by varymid?)? Then at any point on
  interval [a,b], varying slope should give linear increase then decrease? No,
  slope could go to infinity but val's shouldn't. Work it out. Looks like two
  connected concave sections that meet at discontinuous-slope sharp max; not
  concave overall. Maybe not useful.

\item

  Just bandwidth selection using variational-ish approach with exponential
  kernel (vs. exponential distribution over bandwidths)?

\item

  Ensemble of convex kernels using variational-ish approach?

\item

  Combined variational-ish exp-distribution bandwidth selection + kernel
  ensemble / other shaping?

\item

  Choosing kernel bandwidth distribution from some friendly family and
  marginalizing out to get kernel shape?

\item

  Choosing kernel shape in frequency domain? Or various (other) generating functions?

\item

  Optimizing in frequency domain?

\item

  Kernel ``shaping'' optimization problem actually to find how much to geopool /
  timepool? Unfortunately this seems like it requires confronting convexity
  issues or doing a 2-D grid / sampling / Latin hypercube if want to balance two
  factors.  (Although, was there something about log posynomials and
  convexity/concavity dependent on roots or some other criterion?  Might have
  something to check there, but logsumexp GP seems wrong direction.)

\end{itemize}

\subsection{Kernel conditional density estimation}

\subsection{Weighted quantile regression}

\subsection{Multivariate considerations}

\subsection{Localized kernel shapes}

\subsection{Papers to read}
\begin{itemize}
\item \url{https://papers.nips.cc/paper/3393-bayesian-kernel-shaping-for-learning-control}
\item CAKE paper
\item \ldots
\end{itemize}



\end{document}
